\begin{conf-abstract}
{Year-Long miniRUEDI Gas Time Series Reveal Shutdown-Driven Excess Air and Water-Quality Impacts}
{Jared van Rooyen, Yama Tomonaga, Matthias S. Brennwald, Rolf Kipfer, Oliver S. Schilling}
{Eawag}
{Managed aquifer recharge (MAR) operations can unintentionally modulate subsurface gas dynamics, with consequences for redox conditions and groundwater quality. In this study, we deploy a portable mass spectrometer (miniRUEDI) in an observation well at the Hardwald MAR scheme near Basel, Switzerland, and acquire continuous noble- and reactive-gas records (He, Ar, Kr, \ce{N2}, \ce{O2}, \ce{CO2}) over a one-year period. The instrument enabled unattended, on-site measurements with percent-level precision and high temporal resolution, well-suited to capture transient events. During the study, MAR operations ceased eight times for periods ranging from hours to several days. Each cease produced a potential characteristic multiphase gas response: (i) rapid shifts in dissolved gas ratios consistent with bubble entrapment and dissolution upon re-wetting, i.e., formation of “excess air”; (ii) step-changes in \ce{O2} and \ce{N2} that propagated on operational time scales; and (iii) knock-on effects on redox-sensitive species (\ce{CO2}, \ce{Ch4}) indicative of short-term biogeochemical reorganization. By combining noble-gas diagnostics of excess air with reactive-gas trends, we separate physical recharge effects from biogeochemical responses and quantify how shutdowns alter effective recharge signatures and residence-time proxies. These findings highlight portable mass spectrometry as a process-control tool for MAR, linking operational decisions to subsurface gas dynamics and water-quality outcomes in near-real time. Our results generalize established concepts of excess-air formation in porous media to the operational rhythms of MAR and provide actionable guidance for minimizing adverse water-quality swings during start/stop cycles.}
\end{conf-abstract}
